\exo \textit{[Raisonner]}

Un professeur a demandé à ses élèves de déterminer les probabilités des issues d'une expérience aléatoire.

Voici les réponses de Mathilde et Vincent :

Mathilde :

\medskip

\begin{tabular}{*{5}{|c}|}
\hline
Issue &  A & B & C & D  \tabularnewline
\hline
Probabilité & 0,3 & 0,1 & 0,5 & 0,3  \tabularnewline
\hline
\end{tabular}

\medskip

Vincent :

\medskip

\begin{tabular}{*{5}{|c}|}
\hline
Issue &  A & B & C & D  \tabularnewline
\hline
Probabilité & 0,25 & 0,15 & 0,55 & 0,15 \tabularnewline
\hline
\end{tabular}

\medskip

Enzo dit: \og Ils se sont trompés tous les deux \fg{}. Comment le sait-il?

\exo \textit{[Raisonner]}

Le bulletin météo prévoit que, sur une zone donnée, la probabilité de pluie pour demain est 0,4.

Laquelle des affirmations suivantes est la meilleure interprétation du bulletin ?

\begin{enumerate} [label=\Alph*:]
\item Il va pleuvoir sur 40~\% de la zone concernée.
\item Dans cette zone, 40 personnes sur 100 auront de la pluie.
\item Si la même prévision était faite sur 100 jours, il pleuvrait à peu près 40 jours sur 100.
\item Demain, il va pleuvoir pendant 40~\% du temps.
\end{enumerate}

\exo \textit{[Modéliser, calculer]}

\begin{enumerate}
\item Faire l'exercice du manuel scolaire en ligne: LS17 p 309.

\item On suppose que chaque boule est indiscernable au toucher.

Pour chacun des univers vus dans l'exercice, déterminer la loi de probabilité associée.

\end{enumerate}

\exo \textit{[Représenter, Calculer]}

\begin{minipage}{0.7\linewidth}
On fait tourner la roue ci-contre découpée en 11 secteurs identiques et on note la lettre obtenue. On suppose que chaque secteur a la même probabilité de sortir.

Pour chacun des événements suivants, donner les issues qui le réalisent et déterminer sa probabilité :

\begin{itemize}[label=\textbullet]
\item A : \og Obtenir une lettre située avant f dans l'ordre alphabétique \fg{}. 
\item B : \og Obtenir une consonne \fg{}.
\item C : \og Obtenir une lettre composant le mot MATH \fg{}.  
\end{itemize}
\end{minipage}
\hfill
\begin{minipage}{0.3\linewidth}
\begin{center}
\begin{center}
\SpecialCoor
\psset{unit=0.62cm}
\begin{pspicture}(-2.2,-2.35)(2.2,2.35)
  \pscircle[fillstyle=solid,fillcolor=yellow!25,linewidth=0.7pt](0,0){2}
  \multido{\r=90+-32.727}{11}{\psline[linewidth=0.6pt](0,0)(2;\r)}
  \pscircle[fillstyle=solid,fillcolor=white,linewidth=0.5pt](0,0){0.13}
  \rput(1.35;73.6){k}\rput(1.35;40.9){a}\rput(1.35;8.2){b}
  \rput(1.35;-24.5){c}\rput(1.35;-57.3){d}\rput(1.35;-90){e}
  \rput(1.35;-122.7){f}\rput(1.35;-155.5){g}\rput(1.35;171.8){h}
  \rput(1.35;139.1){i}\rput(1.35;106.4){j}
  \pspolygon[fillstyle=solid,fillcolor=black](0,2.18)(-0.1,2.55)(0.1,2.55)
\end{pspicture}
\end{center}
\end{center}
\end{minipage}


\exo \textit{[Raisonner, Calculer]}

Un dé cubique est tel que les probabilités d'apparition des numéros portés par les faces sont données par le tableau suivant, hélas incomplet :

\medskip
\begin{center}
\begin{tabular}{|c|*{6}{p{1 cm}|}}
\hline
Numéro & \centering 1 & \centering 2 & \centering 3 & \centering 4 & \centering 5 & \centering 6  \tabularnewline
\hline
Probabilité & \centering 0,10 & \centering 0,15 & \centering 0,20 & \centering 0,20 & \centering 0,15 &  \tabularnewline
\hline
\end{tabular}
\end{center}


\begin{enumerate}
\item Compléter le tableau.

\item Pour chacun des événements suivants, donner les issues qui le réalisent et déterminer sa probabilité :

\begin{itemize}[label=\textbullet]
\item A : \og Le nombre est pair \fg{}. 
\item B : \og Le nombre est un diviseur de 12 \fg{}.
\item C : \og Le nombre est un premier \fg{}.  
\end{itemize}
\end{enumerate}

\medskip

\exo \textit{[Modéliser, Calculer]}

Deux sélectionneurs Didier et Laurent s'apprêtent à donner leur avis sur la sélection d'un joueur.

On note D et L, respectivement, les événements : \og Didier est pour \fg{} et  \og Laurent est pour \fg{}.

\begin{enumerate}
\item Décrire, dans le contexte, les événements suivants :

a)~$\overline{D}$ \hspace{1cm} b)~$D \cap L$ \hspace{1cm} c)~$D \cup L$ \hspace{1cm} d)~$\overline{D} \cap L$ \hspace{1cm} e)~$\overline{D} \cap \overline{L}$.
\item Que peut-on dire de $D \cup L$ et de $\overline{D} \cap \overline{L}$?
\item On donne: $P(D \cap L)=0,6$ ; $P(D)=0,8$ et $P(L)=0,7$.

Calculer $P(D \cup L)$ et $P(\overline{D} \cap \overline{L})$.
\end{enumerate}

\exo \textit{[Modéliser, Calculer]}

Ce tableau donne la répartition des adhérents à un club de volley-ball.

\medskip

\begin{center}
\renewcommand{\arraystretch}{1.2}
\begin{tabular}{|p{2.5cm}|p{2.5cm}|p{2.5cm}|p{2.5cm}|}
\cline{2-4}
 \multicolumn{1}{c|}{} &  \centering Homme & \centering Femme & \centering Total \tabularnewline
\hline
\centering Jeune &  \centering 27 & \centering 14 & \centering 41 \tabularnewline
\hline
\centering Adulte &  \centering 30 & \centering 29 & \centering 59 \tabularnewline
\hline
\centering  Total &  \centering 57 & \centering 43 & \centering 100 \tabularnewline
\hline
\end{tabular}
\end{center}

On choisit au hasard la licence de l'un des adhérents et on note son sexe et sa catégorie.

\begin{enumerate}
\item Quelle est la probabilité d'avoir choisi :
\begin{enumerate} [label=\alph*)]
\item un jeune ?
\item une femme ?
\item une femme jeune ?
\end{enumerate}
\item Calculer la probabilité d'avoir choisi une femme ou un jeune de deux façons différentes.
\item En déduire la probabilité d'avoir choisi ni une femme ni un jeune.
\end{enumerate}


\exo \textit{[Représenter, Calculer]}

Un village compte 2400 habitants dont 1260 hommes.

\begin{itemize}
\item Un tiers des habitants ont moins de 20 ans.

\item 20~\% des habitants sont des hommes de moins de 20 ans.

\end{itemize}


\medskip

On s'intéresse au prochain habitant qui entrera dans la mairie (qui peut être un enfant suivi de l'un de ses parents). 

\medskip

On considère les événements :

A:\og L'habitant est un homme \fg{};

B: \og L'habitant a plus de 20 ans \fg{}.

\begin{enumerate}
\item Représenter la situation en utilisant un tableau.
\item Déterminer la probabilité de chacun des événements :

a)~$\overline{A}$ \hspace{1cm} b)~$A \cap B$ \hspace{1cm} c)~$\overline{A} \cap B$ 
\item Calculer $P(A \cup B)$ de deux façons différentes.
\item Calculer la probabilité que la personne entrant soit une femme ou soit âgée de plus de 20 ans.
\end{enumerate}  

\exo \textit{[Représenter, Calculer]}

\begin{enumerate}
\item Faire l'exercice CE 5 p 111.
\item Calculer la probabilité des événements A et B.
\item Calculer la probabilité de l'événement $A \cup B$. Interpréter ce résultat.
\end{enumerate}






